\documentclass[11pt]{article}
\usepackage[margin=1in]{geometry}
\usepackage{hyperref}
\usepackage{listings}
\usepackage{xcolor}
\usepackage{enumitem}

\lstset{
    basicstyle=\ttfamily\small,
    backgroundcolor=\color{gray!10},
    frame=single,
    framerule=0.5pt,
    breaklines=true,
    showstringspaces=false,
    language=bash,
    morekeywords={pkg,cmake,ninja,git,cd,mkdir,export,setenv,service,sysrc,sudo,su,doas}
}

\title{Building femm-qt on FreeBSD}
\author{D.\ Meeker}
\date{\today}

\begin{document}
\maketitle

\section{Overview}

femm-qt is a cross-platform port of FEMM (Finite Element Method Magnetics)
using Qt~6 for the GUI and CMake/Ninja for the build system.  This document
describes how to build it from source on FreeBSD.  The instructions target
FreeBSD~14, but should work with minor changes on 13.x and later.

The C++ source and the CMake build are platform-generic: there is no
Linux-specific code, and the only places that differ from the Linux build are
the toolchain (FreeBSD's base compiler is clang, not g++), the package manager
(\texttt{pkg}), and the install prefix (\texttt{/usr/local}, a.k.a.\
\texttt{LOCALBASE}).  The dedicated \texttt{freebsd\_build.sh} script handles
those differences for you.

The build produces six executables:
\begin{itemize}[nosep]
  \item \texttt{femm\_gui} --- the main GUI application
  \item \texttt{femm\_headless} --- headless IPC server (no GUI)
  \item \texttt{fkn} --- magnetics solver
  \item \texttt{belasolv} --- electrostatics solver
  \item \texttt{hsolv} --- heat flow solver
  \item \texttt{csolv} --- current flow solver
\end{itemize}

The commands below use \texttt{sudo} for administrative actions.  GhostBSD
ships \texttt{sudo} preconfigured; on a base FreeBSD system, either install it
(\texttt{pkg install sudo}, run as root) and add your user to the
\texttt{wheel} group, or simply run the privileged commands from a root shell
(\texttt{su -}).  The base system's \texttt{doas} also works once you create
\texttt{/usr/local/etc/doas.conf}.

\section{Prerequisites}

\subsection{Build Tools}

FreeBSD ships clang, the C++ standard library, and the linker in the base
system, so no separate ``build-essential'' package is required.  Install
CMake, Ninja, git, and bash (the build script is a bash script):
\begin{lstlisting}
sudo pkg install cmake ninja git bash
\end{lstlisting}
This provides \texttt{cmake}, \texttt{ninja}, \texttt{git}, and \texttt{bash};
the C++ compiler is already present as \texttt{cc}/\texttt{c++}.

\subsection{Qt 6}

Install the Qt~6 libraries and development files:
\begin{lstlisting}
sudo pkg install qt6
\end{lstlisting}
The \texttt{qt6} metapackage pulls in everything, including the \texttt{moc},
\texttt{rcc}, and \texttt{uic} code generators that the build needs.  If you
prefer a minimal install, \texttt{qt6-base} is usually sufficient for this
project.  Qt~6 installs under \texttt{/usr/local}; no separate Qt installer or
online account is needed.

To verify the installation:
\begin{lstlisting}
qmake6 --version
\end{lstlisting}

\section{Getting the Source}

If you have a source zip, unzip it:
\begin{lstlisting}
unzip femm-qt-src.zip
cd femm-qt
\end{lstlisting}

Or, if the source is hosted on GitHub, clone the repository:
\begin{lstlisting}
git clone <repository-url> femm-qt
cd femm-qt
\end{lstlisting}

\section{Building}

From the \texttt{femm-qt} directory, run the build script:
\begin{lstlisting}
./freebsd_build.sh
\end{lstlisting}

This checks that the build tools are present, configures the project with
CMake (pointing it at \texttt{/usr/local} so Qt~6 is found), and builds it
with Ninja.  All executables are placed in \texttt{build/bin/}.

To wipe the build directory and rebuild from scratch, pass \texttt{clean}:
\begin{lstlisting}
./freebsd_build.sh clean
\end{lstlisting}

If your Qt~6 packages live somewhere other than \texttt{/usr/local}, set
\texttt{LOCALBASE} before running the script:
\begin{lstlisting}
LOCALBASE=/opt/local ./freebsd_build.sh
\end{lstlisting}

If the build fails because CMake cannot find Qt~6, see \textit{Manual
configuration} under Troubleshooting.

\section{Installing}

You can run femm-qt from \texttt{build/bin} indefinitely; nothing requires
an install.  To put it somewhere permanent instead:
\begin{lstlisting}
./freebsd_build.sh install
\end{lstlisting}

This installs what is already built --- it does not rebuild --- into your
home directory:
\begin{lstlisting}
~/.local/opt/femm-qt/   the program, its solvers, the material
                        libraries, the bindings, and a pristine
                        set of the examples
~/.local/bin/           femm_gui, femm_headless, the four solvers
                        (fkn, belasolv, hsolv, csolv), lua4, and
                        femm-lua
~/.local/share/         menu entry and icon
~/Documents/FEMM/
       examples/        your copy of the examples
       docs/            the manual and the scripting references
\end{lstlisting}

The solvers are linked as commands of their own because each is usable
that way: \texttt{fkn model} solves \texttt{model.fem} and writes
\texttt{model.ans}, with \texttt{belasolv}, \texttt{hsolv} and
\texttt{csolv} doing the same for electrostatics, heat flow and current
flow --- which is what you want from a Makefile or a batch queue.
\texttt{lua4} is a stock Lua 4 interpreter and runs directly too.

The examples are installed twice on purpose.  The program directory is out
of the way --- and on Ubuntu worse than that, since the file manager hides
\texttt{\textasciitilde/.local} outright --- so a second copy goes somewhere
you can actually open, under whatever your desktop calls \emph{Documents}.
The shipped examples in it are replaced on every install, so that a
corrected example reaches you rather than only new users --- and because
\emph{running} an example rewrites its \texttt{.fem} (analysis saves the
document first), which makes any ``has this been edited?'' test useless
here.  Anything else you keep in that directory is untouched: only the
files femm-qt ships are written.  Work in a copy, which is what every
example script does on its second line.

FEMM then appears in the application menu.  For \texttt{femm\_gui} to work
as a command, \texttt{\textasciitilde/.local/bin} has to be on your
\texttt{PATH}, and \textbf{the installer does not put it there without being
asked}.  Your shell startup files are your configuration --- kept in a
dotfiles repository, synced between machines, or managed for you on an
institutional login --- and an installer editing one uninvited produces a
confusing diff at best and a change that is silently reverted at worst.

So it checks, and tells you.  If the directory is already on \texttt{PATH}
it says nothing at all.  On Debian and Ubuntu the stock
\texttt{\textasciitilde/.profile} already adds
\texttt{\textasciitilde/.local/bin} when it exists, so it will be there at
your next login and the installer leaves well alone.  Otherwise it prints
the one line to add and which file to add it to.

If you would rather it did the edit, ask for it:
\begin{lstlisting}
./freebsd_build.sh install --modify-path
\end{lstlisting}
which appends a single marked block, once --- running it again finds the
block and adds nothing.

Which file that is depends on your shell, and so does when the change takes
effect.  \texttt{\textasciitilde/.profile} is read at login by the desktop
session, and every terminal inherits the environment from there, so a bash
or \texttt{sh} user must log out and back in --- opening a new terminal
window is \emph{not} enough.  \texttt{\textasciitilde/.zshrc} and fish's
\texttt{config.fish} are read by every interactive shell, so there a new
terminal really is sufficient.  Under \texttt{csh} or \texttt{tcsh} nothing
is written at all: those read none of these files and use a different syntax
entirely, so the installer prints the \texttt{set path} line for
\texttt{\textasciitilde/.cshrc} rather than leave a line that would never run.

Whatever you do, it cannot fix the shell you ran it from --- no program can
change the environment of the shell that launched it.

Pass a different prefix if you would rather it went elsewhere:
\begin{lstlisting}
./freebsd_build.sh install /opt/femm
\end{lstlisting}

It installs into your home rather than \texttt{/usr/local} on purpose.
FEMM finds its solvers, scripting assets and material libraries beside
its own executable, so a single self-contained directory needs no
special handling --- and, more importantly, the material libraries and
the examples stay \emph{writable}.  Both need to be: the material editor
saves the libraries, and every example script opens a model and then
immediately saves it under another name.  In a system directory they
would be read-only and neither would work.

\subsection*{Uninstalling}

\begin{lstlisting}
./freebsd_build.sh uninstall
\end{lstlisting}

It removes the program directory, the commands, the menu entry, the icon
and the documents, and takes out the \texttt{PATH} line if it was the one
that added it.  Pass the same prefix you installed with if it was not the
default, and \texttt{-y} to skip the confirmation.

It works from a manifest the install leaves in the program directory
rather than from a list of paths worked out afresh, because several of
those paths are not reliably re-derivable later --- which solver
commands were linked, whether a profile was edited, and above all where
the documents went, since that comes from \texttt{xdg-user-dir} and can
be reconfigured.  A program directory with no manifest is refused rather
than guessed at, and a command in \texttt{\textasciitilde/.local/bin}
that does not point into the program directory is left alone: a name
like \texttt{lua4} could easily belong to something else.

Your copy of the examples is never removed.  Those are yours, an install
never overwrites them, and a script that deletes a directory somebody has
been working in is not one to run casually.

\section{Running}

Launch the GUI:
\begin{lstlisting}
./build/bin/femm_gui
\end{lstlisting}

Or the headless server:
\begin{lstlisting}
./build/bin/femm_headless
\end{lstlisting}


If you installed as above, use the copies on your \texttt{PATH} instead,
from any directory:
\begin{lstlisting}
femm_gui
femm_headless
\end{lstlisting}

\subsection{Setting FEMM\_DIR}

The language bindings (OctaveFEMM, MathFEMM, pyFEMM, JuliaFEMM) locate the
FEMM executables via the \texttt{FEMM\_DIR} environment variable.  Set it to
the \texttt{bin} directory.  For \texttt{sh}/\texttt{bash} (in
\texttt{\textasciitilde/.profile} or \texttt{\textasciitilde/.bashrc}):
\begin{lstlisting}
export FEMM_DIR=/path/to/femm-qt/build/bin
\end{lstlisting}
For \texttt{csh}/\texttt{tcsh} (the default root shell; in
\texttt{\textasciitilde/.cshrc}):
\begin{lstlisting}
setenv FEMM_DIR /path/to/femm-qt/build/bin
\end{lstlisting}
\textbf{If you installed, you can skip all of this.}  Each binding looks
for \texttt{femm\_gui} and \texttt{femm\_headless} in three places, in
order: \texttt{FEMM\_DIR}, then beside its own directory, then \emph{every
directory on \texttt{PATH}} --- and the install puts them on
\texttt{PATH}.  So the bindings work with nothing set at all.
\texttt{FEMM\_DIR} is how you override that: when running from a build
tree that was never installed, or when you have several builds and need to
say which.  Being first in the list, it beats both the relative search and
\texttt{PATH}.

The binding files themselves are in \texttt{bindings/} under the program
directory --- \texttt{mfiles} for OctaveFEMM, and \texttt{mathfemm},
\texttt{scifemm}, \texttt{juliafemm}, \texttt{pyfemm} and
\texttt{luafemm} beside it.

\subsection{Desktop Integration (Optional)}

\emph{The \texttt{install} step above does this for you, with an icon.}
The recipe below is for running in place out of \texttt{build/bin}.

If you run a desktop environment (KDE, XFCE, etc.), add femm-qt to your
application menu by creating a \texttt{.desktop} file:
\begin{lstlisting}
cat > ~/.local/share/applications/femm-qt.desktop << 'EOF'
[Desktop Entry]
Name=FEMM
Comment=Finite Element Method Magnetics
Exec=/path/to/femm-qt/build/bin/femm_gui
Type=Application
Categories=Science;Engineering;
EOF
\end{lstlisting}

\section{Using with Octave}

MATLAB is not available natively on FreeBSD, so Octave is the supported route
for the *.m language bindings.

Install Octave and a JDK (OctaveFEMM uses Java for socket communication):
\begin{lstlisting}
sudo pkg install octave openjdk21
\end{lstlisting}
Find the path to Java:
\begin{lstlisting}
readlink -f $(which java)
\end{lstlisting}
This usually points under \texttt{/usr/local/openjdk21}.  Set
\texttt{JAVA\_HOME} and \texttt{FEMM\_DIR} in your shell startup file ---
for \texttt{sh}/\texttt{bash}:
\begin{lstlisting}
export JAVA_HOME=/usr/local/openjdk21
\end{lstlisting}
(\texttt{FEMM\_DIR} used to be needed here too.  It no longer is:
\texttt{openfemm} searches \texttt{PATH} for \texttt{femm\_gui} and
\texttt{femm\_headless}, and an install puts them there.  Set it only to
choose between several builds.)
Run Octave and confirm Java is available:
\begin{lstlisting}
version -java
\end{lstlisting}
Then add the OctaveFEMM mfiles to Octave's path:
\begin{lstlisting}
addpath("~/.local/opt/femm-qt/bindings/mfiles")
savepath
\end{lstlisting}
or, working from the source tree instead of an install,
\verb+addpath("/path/to/femm-qt/octavefemm")+.

This is a \emph{different} search path from the one above, and the two are
easy to confuse.  Octave's path is how Octave finds \texttt{openfemm.m};
the system \texttt{PATH} is how \texttt{openfemm.m} then finds
\texttt{femm\_gui}.  Only the second is automatic.  The install prints the
\texttt{addpath} line for every language when it finishes.

A fuller OctaveFEMM guide is installed at
\texttt{\textasciitilde/Documents/FEMM/docs/octavefemm.md}.

\section{Using with Python}

The Python interface is/will be available as pyfemm-qt.  Install Python and
pip first:
\begin{lstlisting}
sudo pkg install python3 py311-pip
\end{lstlisting}
Then install pyfemm from PyPI:
\begin{lstlisting}
pip install pyfemm-qt
\end{lstlisting}
To install the version that shipped with this source instead, change to the
\texttt{pyfemm} directory --- in the source tree, or
\verb+~/.local/opt/femm-qt/bindings/pyfemm+ after an install --- and install
in editable mode:
\begin{lstlisting}
cd pyfemm
pip install -e .
\end{lstlisting}
\texttt{-e} records the \emph{absolute path} of that directory, so an
editable install breaks if the directory is later moved or deleted.  Point
it at the installed copy, not at a build tree you might clean.
Like recent Linux distributions, FreeBSD marks the system Python as
externally managed (PEP~668), so pip may refuse to install into it.  The
cleanest fix is a virtual environment:
\begin{lstlisting}
python3 -m venv ~/femm-venv
. ~/femm-venv/bin/activate
pip install -e .
\end{lstlisting}
Since pyfemm is a single file with no dependencies, overriding the check is
also safe:
\begin{lstlisting}
pip install -e . --break-system-packages
\end{lstlisting}
A fuller pyFEMM guide --- multi-instance use, headless mode, connecting to
an instance that is already running --- is installed at
\texttt{\textasciitilde/Documents/FEMM/docs/pyfemm.md}.

\section{Troubleshooting}

\subsection{Manual configuration}

The build script just wraps the standard CMake and Ninja steps.  To run them
by hand:
\begin{lstlisting}
mkdir build && cd build
cmake .. -G Ninja -DCMAKE_BUILD_TYPE=Release \
    -DCMAKE_PREFIX_PATH=/usr/local
ninja
\end{lstlisting}
The \texttt{CMAKE\_PREFIX\_PATH=/usr/local} hint lets find\_package locate the
Qt~6 packages installed by \texttt{pkg}.

\subsection{``env: bash: No such file or directory''}

The build script is a bash script, but bash is not part of the FreeBSD base
system.  Install it:
\begin{lstlisting}
sudo pkg install bash
\end{lstlisting}

\subsection{Qt not found}

If CMake reports that it cannot find Qt6, check that the package is installed:
\begin{lstlisting}
pkg info -x qt6
\end{lstlisting}
If it is installed but not found, pass \texttt{-DCMAKE\_PREFIX\_PATH} pointing
at the Qt~6 CMake config directory, typically
\texttt{/usr/local/lib/cmake/Qt6}:
\begin{lstlisting}
cmake .. -G Ninja -DCMAKE_BUILD_TYPE=Release \
    -DCMAKE_PREFIX_PATH=/usr/local/lib/cmake/Qt6
\end{lstlisting}

\subsection{Missing OpenGL headers}

On minimal installations or VMs without a desktop, Qt may need the Mesa
OpenGL libraries and DRI drivers:
\begin{lstlisting}
sudo pkg install mesa-libs mesa-dri libglvnd
\end{lstlisting}

\subsection{VirtualBox Guest Additions}

If running in a VirtualBox VM, install the FreeBSD guest additions for proper
display scaling and shared folders:
\begin{lstlisting}
sudo pkg install virtualbox-ose-additions
\end{lstlisting}
Enable the guest services so they start at boot:
\begin{lstlisting}
sudo sysrc vboxguest_enable="YES"
sudo sysrc vboxservice_enable="YES"
\end{lstlisting}
Reboot after installation.  To use a shared folder, add your user to the
\texttt{wheel}/\texttt{operator} group as appropriate and mount it with the
\texttt{vboxsf} filesystem, e.g.:
\begin{lstlisting}
sudo mount -t vboxvfs <share-name> /mnt
\end{lstlisting}

\subsection{MESA / OpenGL errors in VirtualBox}

VirtualBox's virtual GPU may not support the OpenGL version Qt expects,
producing errors like ``ZINK: failed to choose pdev'' or ``failed to create
dri2 screen.''  Two fixes:

\begin{enumerate}[nosep]
\item Enable 3D Acceleration in VirtualBox (Settings $\rightarrow$ Display
      $\rightarrow$ check ``Enable 3D Acceleration'' and raise the video
      memory to at least 128~MB in Expert mode), or
\item Force Qt to use software rendering by adding (for \texttt{sh}/\texttt{bash})
\begin{lstlisting}
export LIBGL_ALWAYS_SOFTWARE=1
\end{lstlisting}
to your shell startup file (or \texttt{setenv LIBGL\_ALWAYS\_SOFTWARE 1} for
csh/tcsh).
\end{enumerate}

\end{document}
